\section{树链剖分}

\par \noindent \textbf{问题场景：} 链上求和，链上求最值，链上修改，子树修改，子树求和。
~\\
\par \noindent 如：修改和查询点 x 到 y 的路径的信息；修改和查询以 x 为根的子树的信息。

\subsection{树链剖分（路径与子树操作）}

\par \noindent 第一次 DFS 求出每个点的父亲、深度、子树大小和重儿子；第二次 DFS 按重儿子优先的顺序给每个点分配新编号，使同一条重链上的点编号连续，且每棵子树的编号也连续。

\begin{minted}{c++}
const int N=2e5+10;
struct qwq
{
    int to,ne;
}e[N];
int n,m,root,mod;
int num,elast[N];
void add(int x,int y)
{
    e[++num]={y,elast[x]},elast[x]=num;
}

struct xds
{
    int l,r,sum,lazy;
}t[N<<2];
void build(int p,int l,int r)
{
    t[p].l=l,t[p].r=r;
    if(l==r)return;
    int mid=(l+r)>>1;
    build(p*2,l,mid),build(p*2+1,mid+1,r);
}
void pushdown(int p)
{
    t[p*2].lazy=(t[p*2].lazy+t[p].lazy)%mod;
    t[p*2+1].lazy=(t[p*2+1].lazy+t[p].lazy)%mod;
    t[p*2].sum=(t[p*2].sum+(t[p*2].r-t[p*2].l+1)*t[p].lazy%mod)%mod;
    t[p*2+1].sum=(t[p*2+1].sum+(t[p*2+1].r-t[p*2+1].l+1)*t[p].lazy%mod)%mod;
    t[p].lazy=0;
}
void modify(int p,int l,int r,int x)
{
    if(l<=t[p].l&&t[p].r<=r)
    {
        t[p].lazy=(t[p].lazy+x)%mod;
        t[p].sum=(t[p].sum+(t[p].r-t[p].l+1)*x%mod)%mod;
        return;
    }
    pushdown(p);
    int mid=(t[p].l+t[p].r)>>1;
    if(l<=mid)modify(p*2,l,r,x);
    if(r>=mid+1)modify(p*2+1,l,r,x);
    t[p].sum=(t[p*2].sum+t[p*2+1].sum)%mod;
}
int query(int p,int l,int r)
{
    if(l<=t[p].l&&t[p].r<=r)return t[p].sum;
    pushdown(p);
    int mid=(t[p].l+t[p].r)>>1,res=0;
    if(l<=mid)res=(res+query(p*2,l,r))%mod;
    if(r>=mid+1)res=(res+query(p*2+1,l,r))%mod;
    return res;
}

int fa[N],dep[N],si[N],son[N],top[N],dfn[N],tot;
void dfs1(int x,int u)
{
    dep[x]=dep[u]+1,fa[x]=u;
    for(int i=elast[x];i;i=e[i].ne)
    {
        int to=e[i].to;
        if(to==u)continue;
        dfs1(to,x);
        si[x]+=si[to];
        if(si[to]>si[son[x]])son[x]=to;
    }
    si[x]++;
}
void dfs2(int x,int u)
{
    top[x]=u,dfn[x]=++tot;
    if(!son[x])return;
    dfs2(son[x],u);
    for(int i=elast[x];i;i=e[i].ne)
    {
        int to=e[i].to;
        if(to!=son[x]&&to!=fa[x])dfs2(to,to);// ***
    }
}
int ask(int x,int y)
{
    int res=0;
    while(top[x]!=top[y])
    {
        if(dep[top[x]]<dep[top[y]])swap(x,y);
        res=(res+query(1,dfn[top[x]],dfn[x]))%mod;
        x=fa[top[x]];
    }
    if(dep[x]>dep[y])swap(x,y);
    res=(res+query(1,dfn[x],dfn[y]))%mod;
    return res;
}
void change(int x,int y,int k)
{
    while(top[x]!=top[y])
    {
        if(dep[top[x]]<dep[top[y]])swap(x,y);
        modify(1,dfn[top[x]],dfn[x],k);
        x=fa[top[x]];
    }
    if(dep[x]>dep[y])swap(x,y);
    modify(1,dfn[x],dfn[y],k);
}

int w[N];

signed main()
{
    n=read(),m=read(),root=read(),mod=read();
    for(int i=1;i<=n;i++)w[i]=read();
    
    for(int i=1;i<n;i++)
    {
        int x=read(),y=read();
        add(x,y),add(y,x);
    }
    dfs1(root,0),dfs2(root,root);
    build(1,1,n);
    
    for(int i=1;i<=n;i++)modify(1,dfn[i],dfn[i],w[i]);
    
    for(int i=1;i<=m;i++)
    {
        int op=read();
        if(op==1)
        {
            int x=read(),y=read(),z=read()%mod;
            change(x,y,z);
        }
        if(op==2)
        {
            int x=read(),y=read();
            cout<<ask(x,y)<<endl;
        }
        if(op==3)
        {
            int x=read(),k=read()%mod;
            modify(1,dfn[x],dfn[x]+si[x]-1,k);
        }
        if(op==4)
        {
            int x=read();
            cout<<query(1,dfn[x],dfn[x]+si[x]-1)<<endl;
        }
    }
    return 0;
}
\end{minted}

\par \noindent 说明：
\begin{itemize}
\item \textbf{dfs1}：计算每个点的深度、父亲、子树大小和重儿子，为后续剖分做准备。
\item \textbf{dfs2}：沿重儿子优先分配 DFS 序，使同一条重链上节点的 \texttt{dfn} 连续，且任意子树内节点的 \texttt{dfn} 也连续。因此子树的区间为 \texttt{[dfn[x], dfn[x]+si[x]-1]}。
\item \textbf{ask(x,y)}：路径查询。每次让所在链顶深度较大的点跳到父节点，直到两点在同一重链上，最后在同一链上查询连续区间。
\item \textbf{change(x,y,k)}：路径修改，逻辑与查询相同，只是把区间查询换成区间加。
\item \textbf{子树操作}：直接对区间 \texttt{[dfn[x], dfn[x]+si[x]-1]} 做区间加或区间求和即可。
\end{itemize}

\subsection{树链剖分维护 LCA}

\par \noindent 在 \texttt{dfs1} 和 \texttt{dfs2} 完成之后，可以利用 \texttt{top} 和 \texttt{dep} 在 $O(\log n)$ 时间内求 LCA。

\begin{minted}{c++}
int lca(int x,int y)
{
    while(top[x]!=top[y])
    {
        if(dep[top[x]]<dep[top[y]])
            swap(x,y);
        x=fa[top[x]];
    }
    return dep[x]<dep[y]?x:y;
}
int dis(int x,int y)
{
    return dep[x]+dep[y]-2*dep[lca(x,y)];
}
\end{minted}

\subsection{dfs 序 + 倍增 O(1) 维护 LCA}

\par \noindent 先用 DFS 求出欧拉序（进入节点时记录），再对欧拉序上深度最小的节点建立 ST 表，即可在 $O(1)$ 时间内回答 LCA。

\begin{minted}{c++}
// P3379 【模板】最近公共祖先（LCA）https://www.luogu.com.cn/problem/P3379
const int N=1000010;

int h[N],e[N<<1],ne[N<<1],idx;
void add(int a,int b){e[idx]=b,ne[idx]=h[a],h[a]=idx++;}
int dfn[N],timstamp;
int n,m,rt;
int dep[N],rev[N],fa[N];

void dfs(int u)
{
    dep[u]=dep[fa[u]]+1;
    dfn[u]=++timstamp;
    rev[timstamp]=u;
    for(int i=h[u];i!=-1;i=ne[i])
    {
        int v=e[i];
        if(v==fa[u]) continue;
        fa[v]=u;
        dfs(v);
    }
}
inline int MIN(int x,int y){return dep[x]<dep[y]?x:y;}

int st[N][21];
int lg[N];
int lca(int u,int v)
{
    if(u==v) return u;
    u=dfn[u],v=dfn[v];
    if(u>v) swap(u,v);
    u++;
    int k=lg[v-u+1];
    return fa[MIN(st[u][k],st[v-(1<<k)+1][k])]; 
}
signed main()
{
    ios::sync_with_stdio(false);cin.tie(nullptr);cout.tie(nullptr);
    cin>>n>>m>>rt;
    for(int i=1;i<=n;i++) h[i]=-1;
    for(int i=1;i<n;i++)
    {
        int a,b;cin>>a>>b;
        add(a,b),add(b,a);
    }
    dfs(rt);
    for(int i=2;i<=n;i++) lg[i]=lg[i>>1]+1;
    for(int i=1;i<=n;i++) st[i][0]=rev[i];
    for(int k=1;k<=lg[n];k++) 
        for(int i=1;i+(1<<k)-1<=n;i++) 
            st[i][k]=MIN(st[i][k-1],st[i+(1<<k-1)][k-1]);

    while(m--)
    {
        int x,y;cin>>x>>y;
        cout<<lca(x,y)<<'\n';
    }
    return 0;
}
\end{minted}

\subsection{长链剖分}

\par \noindent 长链剖分的价值主要体现在能优化树上 \textbf{与深度有关} 的 DP。如果子树内 \textbf{每个深度仅有一个信息}，就可以使用长链剖分优化。一般形式如：设 $f(i,j)$ 表示以 $i$ 为根的子树内，深度为 $j$ 的节点的贡献。

~\\
\par \noindent \textbf{经典结论：}选一个节点能覆盖它到根的所有节点。选 $k$ 个节点，覆盖的最多节点数就是前 $k$ 条长链长度之和，选择的节点即 $k$ 条长链末端。一个解决方案是使用指针动态申请内存：对于一条重链，共用一个大小为其长度的数组。另一个方案是使用 vector 的 swap 特性。

~\\
\par \noindent 长链剖分实现起来有很多细节，例如如何 \textbf{继承重儿子} 的 DP 值，以及如何处理合并时 \textbf{下标偏移} 的问题。

\begin{tcolorbox}
\par \noindent 给定一棵以 1 为根，$n$ 个节点的树。设 $d(u,x)$ 为 $u$ 子树中到 $u$ 距离为 $x$ 的节点数。

\par \noindent 对于每个点，求一个最小的 $k$，使得 $d(u,k)$ 最大。
\end{tcolorbox}

\par \noindent 类似启发式合并，每次继承\textbf{"重"儿子}的答案，然后将所有轻儿子的答案合并过来。时间复杂度是优秀的 $O(n)$

\begin{minted}{c++}
// 指针版本
const int N=1000010;
vector<int> e[N];
int n;
int pool[N];
int *f[N],*now=pool;
int ans[N];
int dep[N],son[N];

void dfs1(int u,int fa)
{
    for(int v:e[u])
    {
        if(v==fa) continue;
        dfs1(v,u);
        if(dep[v]>dep[son[u]]) son[u]=v;
    }
    dep[u]=dep[son[u]]+1;
}
void dfs2(int u,int fa)
{
    f[u][0]=1;

    if(son[u])
    {
        f[son[u]]=f[u]+1;
        dfs2(son[u],u);
        ans[u]=ans[son[u]]+1;
    }
    for(int v:e[u])
    {
        if(v==fa||v==son[u]) continue;
        f[v]=now;
        now+=dep[v];
        dfs2(v,u);
        for(int i=1;i<=dep[v];i++)
        {
            f[u][i]+=f[v][i-1];
            if(f[u][i]>f[u][ans[u]]||f[u][i]==f[u][ans[u]]&&i<ans[u])
                ans[u]=i;
        }
    }
    if(f[u][ans[u]]==1) ans[u]=0;
}

signed main()
{
    ios::sync_with_stdio(false);cin.tie(nullptr);cout.tie(nullptr);
    cin>>n;
    for(int i=1,u,v;i<n;i++)
    {
        cin>>u>>v;
        e[u].push_back(v);
        e[v].push_back(u);
    }
    dfs1(1,0);
    f[1]=now;
    now+=dep[1];
    dfs2(1,0);
    for(int i=1;i<=n;i++) cout<<ans[i]<<'\n';
    return 0;
}
\end{minted}

\begin{minted}{c++}
// vector 版本
const int N=1e6+5;
int n,dep[N],mxd[N],son[N],fa[N];
int mx[N],ans[N];
vector<int> e[N],f[N]; // 设 f[i][j] 表示点 i 的 j 级子节点的数量。
void dfs1(int u,int p)
{
    dep[u]=dep[p]+1,fa[u]=p;
    for(int v:e[u])
    {
        if(v==p) continue;
        dfs1(v,u),mxd[u]=max(mxd[u],mxd[v]+1);
        if(mxd[v]>mxd[son[u]]) son[u]=v;
    }
}
// 不妨将信息倒过来存储在 vector 中，转化为在动态数组末端插入一个数。
/**
1. 按深度递增的顺序存储的话，因为合并重儿子信息时要在开头插入元素，效率低下。
   所以考虑按深度递减的顺序存储信息。
2. 合并重儿子信息的时候，直接用 swap 交换而不是复制，在时间和空间上都更优
  （swap 交换 vector 的时间复杂度是 O(1) 的）。
**/
void dfs2(int u)
{
    if(son[u])
    {
        dfs2(son[u]),f[u].swap(f[son[u]]);
        mx[u]=mx[son[u]],ans[u]=ans[son[u]];
    }
    for(int v:e[u])
    {
        if(v==fa[u]||v==son[u]) continue;
        dfs2(v);
        for(int i=0;i<f[v].size();i++)
        {
            int d=f[v].size()-i-1,p=f[u].size()-d-1;
            f[u][p]+=f[v][i];
            if(f[u][p]>mx[u]||f[u][p]==mx[u]&&ans[u]>d) mx[u]=f[u][p],ans[u]=d;
        }
    } 
    f[u].push_back(1);
    if(mx[u]<=1) mx[u]=1,ans[u]=0;
    else ans[u]++;
}
signed main()
{
    cin>>n,mxd[0]=-1;
    for(int i=1;i<n;i++)
    {
        int u,v;cin>>u>>v;
        e[u].push_back(v),e[v].push_back(u);
    } 
    dfs1(1,0),dfs2(1);
    for(int i=1;i<=n;i++) cout<<ans[i]<<'\n';
    return 0;
}
\end{minted}